\chapter{\MakeUppercase{Project Description and Goals}}
\justifying
\section{Literature Review}
\paragraph{} This chapter reviews the current literature related to dynamic difficulty adjustment and reinforcement learning applied to gaming. These works provided the theoretical basis and inspiration for this project, establishing the foundations upon which this work is built.

\subsection{Reinforcement Learning in Games}
\paragraph{} \ac{rl} has a long history in game-playing agents. \citet{mnih2015human} introduced Deep Q-Networks (DQN), showing that \ac{rl} agents could achieve human-level performance in Atari games through the use of experience replay and target networks. This breakthrough established deep \ac{rl} as a workable method for complicated decision-making in interactive ecosystems.

\paragraph{} \citet{silver2016mastering} made a significant leap forward with AlphaGo, combining Monte Carlo Tree Search with deep neural networks to defeat a world champion in the game of Go. Their work showed that \ac{rl} agents could learn to play games that had huge state spaces through self-play and learned evaluation functions.

\subsection{Proximal Policy Optimization}
\paragraph{} \citet{schulman2017proximal} introduced \ac{ppo}, a policy gradient method that alternates between sampling data by interacting with the environment and optimizing a surrogate objective function using stochastic gradient ascent. \ac{ppo} takes care of the instability of standard policy gradient methods by limiting the size of updates through a clipped probability ratio. This feature makes \ac{ppo} especially well suited for environments calling for constant adaptation, such as \ac{dda} systems.
\paragraph{} To fully appreciate the role of \ac{ppo} in this project, one must examine the fundamental mechanics of Policy Gradient methods. Unlike value-based methods such as DQN, which attempt to learn the value of taking a specific action in a given state (the Q-value), Policy Gradient methods directly optimize the parameterized policy function $\pi(a|s, \theta)$, where $\theta$ represents the weights of the neural network. By directly parameterizing the policy, these algorithms are natively capable of handling continuous action spaces and stochastic policies. This is highly advantageous in \ac{dda}, where adjustments often require high degrees of granularity (such as minor, continuous increments in enemy speed).
\paragraph{} The primary objective in standard Policy Gradients is to maximize the expected sum of rewards, denoted mathematically as $J(\theta) = \mathbb{E}_{\pi_\theta} [R(\tau)]$. The gradient of this objective, derived via the Policy Gradient Theorem, is $\nabla_\theta J(\theta) = \mathbb{E}_{\pi_\theta} [\nabla_\theta \log \pi_\theta(a|s) Q^{\pi_\theta}(s,a)]$. However, native policy gradients suffer from high variance and severe sample inefficiency due to massive, destructive step changes during the optimization phase. \ac{ppo} smartly limits this volatility by defining a surrogate objective bound. The objective introduces a probability ratio $r_t(\theta) = \frac{\pi_\theta(a_t|s_t)}{\pi_{\theta_{old}}(a_t|s_t)}$, combined with a clipping mechanism: $L^{CLIP}(\theta) = \mathbb{E}_t [\min(r_t(\theta) \hat{A}_t, \text{clip}(r_t(\theta), 1-\epsilon, 1+\epsilon)\hat{A}_t)]$. This explicit mathematical clipping bounds the policy update interval, entirely preventing catastrophic forgetting and ensuring monotonic improvement during the agent's gameplay tuning phase.

\subsection{Analysis of Player Modeling}
\paragraph{} The effectiveness of \ac{dda} algorithms heavily depends upon the underlying schema governing player behavior analysis. Player modeling in video games revolves around extracting semantic intent and skill classification from raw, abstract telemetry data. Early taxonomies, such as Bartle’s Taxonomy of Player Types, categorized individuals into conceptual archetypes (Achievers, Explorers, Socializers, Killers) based heavily on their interaction motivations. Modern \ac{dda} literature has increasingly attempted to map in-game telemetry directly to these psychological profiles, creating dynamic "Persona-Driven" difficulty systems.
\paragraph{} In high-speed, real-time action games, this profiling must be conducted instantaneously via temporal data streams. The literature frequently identifies survival duration, tracking accuracy, reaction time (latency between stimulus and input), and spatial positioning as core, universal markers of physical skill. For instance, erratic movement vector patterns are strongly correlated with high cognitive load, panic, and frustration, whereas smooth, predictive movement trajectories reliably signal flow state and mastery. Tracking the localized variance of these metrics---rather than solely calculating their mean values---offers a much richer psychological fingerprint of the current player. This project incorporates these sophisticated insights by leveraging a multi-dimensional, continuous observation vector representing game state. This complex array allows the underlying \ac{ppo} agent to implicitly model the latent psychological state of the player without having to explicitly pigeonhole them into rigid, unyielding personality profiles.

\subsection{Dynamic Difficulty Adjustment}
\paragraph{} \citet{hunicke2004ai} developed one of the earliest formalized frameworks for \ac{dda}, introducing the concept of player modeling and adaptive challenge. Their work established the theoretical foundation for difficulty adjustment as a feedback control problem, where the player's performance measures act as inputs and the game parameters are the outputs.

\paragraph{} \citet{andrade2006dynamic} used player behavior analysis to create a \ac{dda} system in which player actions were clustered and skill levels identified, then difficulty was adapted accordingly. Their methodology demonstrated the effectiveness of continuous tracking over checkpoint-based assessment.

\paragraph{} \citet{zohaib2018dynamic} carried out a comprehensive survey of \ac{dda} techniques, categorizing them into rule-based, probabilistic, and machine learning-based approaches. They focused attention on \ac{rl} as a promising direction since it can learn a policy that adapts to hidden dynamics without explicit programming of adjustment rules.

\subsection{RL-Based Dynamic Difficulty Adjustment}
\paragraph{} \citet{zhao2020deep} applied deep \ac{rl} to \ac{dda} in a fighting game context, implementing a DQN agent that adjusted enemy \ac{ai} behavior in response to player performance. Their research demonstrated that \ac{rl}-based \ac{dda} retained higher player engagement compared to static difficulty settings.

\paragraph{} \citet{silva2021dynamic} considered applying policy gradient techniques to \ac{dda} in educational games, showing that \ac{rl} agents could learn to maintain a balance of challenge and learning outcomes simultaneously. Their research demonstrated the potential of \ac{ppo} for real-time difficulty adjustment tasks.

\section{Research Gap}
\paragraph{} Based on the literature review, the following research gaps have been identified:

\begin{itemize}
	\item \textbf{Real-Time Space Shooter DDA:} Most existing \ac{rl}-based \ac{dda} research focuses on turn-based or strategy games. Real-time application of \ac{ppo}-based \ac{dda} to action games, especially space shooters, remains underexplored.
	\item \textbf{Multi-Metric Monitoring:} Most \ac{dda} systems depend upon a single performance metric (for example, score or health). Comprehensive multi-metric monitoring of accuracy, kill rate, health, damage rate, and survival time for difficulty balancing is less studied.
	\item \textbf{Full-Stack Integration:} There is limited work on the integration of \ac{rl}-based \ac{dda} with live monitoring and visualization tools such as analytics dashboards.
	\item \textbf{Temporal Modeling:} The utilization of \ac{lstm} layers to capture temporal trends in player behavior for more anticipatory difficulty adjustments is a novel contribution.
\end{itemize}

\section{Objectives}
\paragraph{} The primary objectives of this project are:

\begin{enumerate}
	\item To develop an adaptive gameplay system using \ac{rl} in a Pygame-based space shooter game.
	\item To implement \ac{ppo} as the core algorithm for dynamic difficulty adjustment, training an agent to vary the enemy speed, health, spawn delay, and attack intensity based on live player performance.
	\item To design a comprehensive player performance monitoring system tracking five key metrics: accuracy ratio, kill rate, player health, damage rate, and survival time.
	\item To build a FastAPI backend for logging gameplay data and a React frontend dashboard for real-time performance visualization.
	\item To evaluate the effectiveness of the \ac{dda} system through comparative playtesting with static difficulty.
	\item To extend the base model with \ac{lstm} layers for temporal player behavior analysis.
\end{enumerate}

\section{Problem Statement}
\paragraph{} Video games with static difficulty settings fail to accommodate the diverse skill levels of players. For better players, the game becomes boring, while for beginners it becomes frustrating, both of which lead to player disengagement. This project attempts to address this problem by developing an \ac{rl}-based \ac{dda} system using \ac{ppo} that automatically adapts game difficulty parameters in real time through continuous player performance monitoring. The system has been implemented as part of a space shooter game with a complete analytics pipeline for monitoring and visualization.

\section{Project Plan}
\paragraph{} The project follows a structured development plan organized into six phases:

\begin{enumerate}
	\item \textbf{Phase 1 -- Game Development (Weeks 1--3):} Develop the Pygame-based space shooter game (\texttt{game\_3d.py}) with player controls, enemy system, asteroid mechanics, and collision detection.
	\item \textbf{Phase 2 -- Environment Wrapping (Weeks 4--5):} Create the \ac{rl} environment wrapper (\texttt{env\_wrapper\_3d.py}) with a Gymnasium-compatible interface that defines observation and action spaces.
	\item \textbf{Phase 3 -- PPO Training (Weeks 6--8):} Implement the \ac{ppo} training pipeline (\texttt{train\_dda\_3d.py}), define the reward function, and train the \ac{dda} agent.
	\item \textbf{Phase 4 -- Integration and Backend (Weeks 9--10):} Integrate the trained model into the game for live inference. Create the backend using FastAPI for data logging.
	\item \textbf{Phase 5 -- Dashboard Development (Weeks 11--12):} Build the React analytics dashboard to track and visualize player data and difficulty trends in real time.
	\item \textbf{Phase 6 -- Testing and Evaluation (Weeks 13--14):} Conduct playtesting, evaluate the \ac{dda} system, implement the \ac{lstm} extension, and finalize documentation.
\end{enumerate}
